\section{Methods}
\label{sec:methods}

\subsection{Dataset}

We generate a synthetic cohort of 200 patients using Synthea \cite{walonoski2018synthea},
an open-source patient simulator that produces medically realistic FHIR R4 records based
on clinical guidelines, public health statistics, and US census data. Each patient is
aged 40--75 at simulation end, has at least 10 years of medication history, and has at
least one currently active medication. Because Synthea creates one \texttt{MedicationRequest}
resource per prescription refill, a patient with a 25-year chronic disease history can
accumulate hundreds of entries for a handful of distinct drugs. No real patient data is
used at any point.

Ground truth is defined as the set of distinct medication names from all
\texttt{MedicationRequest} resources with \texttt{status = "active"} in the FHIR bundle.
Names are taken from the \texttt{medicationCodeableConcept.text} field, exactly as they
appear in the record. The ground truth set for the 200 patients ranges from 1 to 16 active
medications, with a median of 5.

\subsection{Serialisation Strategies}
\label{sec:strategies}

All four strategies receive the same FHIR bundle and produce a text string that is
embedded verbatim into the model prompt. We strip administrative fields
(\texttt{subject}, \texttt{encounter}, \texttt{requester}, \texttt{meta}, \texttt{id})
from all strategies; the model sees only clinically relevant content. Table~\ref{tab:strategy_comparison}
summarises how the four strategies differ along the dimensions that matter for model
comprehension: structural form, ordering, and how the active/historical distinction is
signalled. Appendix~\ref{app:strategy_examples} shows the complete serialised output of
each strategy for a single patient.

\begin{table*}[t]
\centering
\small
\renewcommand{\arraystretch}{1.25}
\begin{tabular}{p{2.4cm} p{3.2cm} p{3.2cm} p{3.2cm} p{3.2cm}}
\toprule
\textbf{Aspect} & \textbf{A: Raw JSON} & \textbf{B: Markdown Table} & \textbf{C: Clinical Narrative} & \textbf{D: Chrono. Timeline} \\
\midrule
Structural form &
Nested FHIR JSON, indented &
Six-column markdown table &
One sentence per medication &
Pipe-delimited timeline lines \\

Ordering &
As in the FHIR bundle, grouped by resource &
Chronological, oldest to newest &
Active medications first, then history &
Strictly chronological, oldest to newest \\

Active/historical signal &
\texttt{"status"} field inside each resource &
\texttt{active} or \texttt{completed} cell in the Status column &
Explicit heading: \textit{``Currently active medications:''} &
Inline \texttt{active} token on each line, no separation \\
\bottomrule
\end{tabular}
\caption{The four serialisation strategies compared along their structural, ordering,
  and status-signalling dimensions. Full verbatim examples of each strategy applied
  to the same patient appear in Appendix~\ref{app:strategy_examples}.}
\label{tab:strategy_comparison}
\end{table*}

\paragraph{Strategy A — Raw JSON (baseline).}
\texttt{MedicationRequest} resources are serialised as indented JSON with clinical
fields intact. Because Synthea emits one resource per refill, the full JSON bundle
for some patients exceeds the context windows of the smaller models in our evaluation
(Mistral-7B and BioMistral-7B at 32K tokens). For Strategy A only, we therefore cap
the input at 100 resources: all active medications are kept, and the remaining slots
are filled with the most recent historical records. Strategies B, C, and D apply no
such cap because their more compact representations (a flat table, per-medication
sentences, and pipe-delimited lines) fit the full history within context for every
patient in the cohort. This asymmetry is itself part of what Strategy A represents
in practice: a raw-JSON pipeline is not deployable for typical chronic-disease
histories without some form of truncation or preprocessing, so the active-first cap
is a minimally-invasive form of the filtering any real deployment would have to
perform. It is a conservative choice for Strategy A in our comparison: active items
are guaranteed to appear, so any remaining performance gap against the other
strategies reflects format effects rather than missing input.

\paragraph{Strategy B — Markdown Table.}
All \texttt{MedicationRequest} resources are flattened into a single markdown table
(columns: Medication, RxNorm code, Status, Prescribed date, Dose, Frequency), sorted
chronologically. Missing dose or frequency fields are shown as a dash. This format is
structurally explicit and human-readable while remaining machine-parseable.

\paragraph{Strategy C — Clinical Narrative.}
Each medication becomes a plain-English sentence. Active medications appear first under
a labelled heading (\textit{Currently active medications:}), followed by a separate
section for historical medications. This structure explicitly signals the task-relevant
partition, reducing the model's need to classify status from raw field values.

\paragraph{Strategy D — Chronological Timeline.}
All \texttt{MedicationRequest} resources are sorted strictly by prescription date,
oldest to newest, regardless of status. Each entry is a pipe-delimited line showing
date, status, medication name, and dosage. This format makes no attempt to distinguish
active from historical medications — the model must reason across the full temporal
sequence to determine what is currently active.

\subsection{Models}
\label{sec:models}

We evaluate five open-weight models spanning two orders of magnitude in parameter count.
Table~\ref{tab:models} summarises specifications. All models are served locally via
Ollama using 4-bit quantisation on an AWS \texttt{g6e.xlarge} instance
(NVIDIA L40S, 48 GB VRAM). No proprietary APIs are used.

\begin{table*}[t]
\centering
\small
\renewcommand{\arraystretch}{1.25}
\begin{tabular}{lrrrll}
\toprule
\textbf{Model} & \textbf{Params} & \textbf{Native Ctx} & \textbf{Runtime Ctx} & \textbf{Type} & \textbf{Quant.} \\
\midrule
Phi-3.5-mini   & 3.8B  & 128K & 64K          & Instruct & Q4\_0   \\
Mistral-7B     & 7B    & 32K  & 64K$^{*}$    & Instruct & Q4\_K\_M \\
BioMistral-7B  & 7B    & 32K  & 64K$^{*}$    & Pretrain & Q4\_K\_M \\
Llama-3.1-8B   & 8B    & 128K & 64K          & Instruct & Q4\_K\_M \\
Llama-3.3-70B  & 70B   & 128K & 32K$^{\dag}$ & Instruct & Q4\_K\_M \\
\bottomrule
\end{tabular}
\caption{Models evaluated. \textit{Native Ctx} is the model's trained maximum context;
  \textit{Runtime Ctx} is the \texttt{num\_ctx} setting passed to Ollama for our runs.
  $^{*}$Mistral-7B and BioMistral-7B were run with a runtime context exceeding their
  native 32K window; this affects 13 strategy-A patients whose serialised prompts
  exceed 32K tokens. $^{\dag}$Llama-3.3-70B was capped at 32K to keep weights plus KV
  cache within the 48 GB L40S VRAM budget. BioMistral is a further-pretrained base
  model (not instruction-tuned), included to test whether domain pretraining alone is
  sufficient. All other models are instruction-tuned.}
\label{tab:models}
\end{table*}

BioMistral~\cite{labrak2024biomistral} is further-pretrained on PubMed Central text but
is \textit{not} instruction-tuned. We include it specifically to test whether domain
pretraining alone is sufficient for a structured extraction task, given that several
recent deployments use completion-style medical models for clinical NLP tasks.

\subsection{Prompt Design}
\label{sec:prompt}

We use a single prompt template across all models and all strategies:

\begin{quote}
\small
\textit{You are a clinical assistant performing medication reconciliation.
You will be given a patient's medication history. Your task is to identify all
medications that are currently ACTIVE for this patient.
A medication is currently active if its status is ``active''. Medications with
status ``completed'', ``stopped'', ``cancelled'', or ``on-hold'' are historical
and must NOT be included in your answer.
Return your answer as a JSON array of medication names exactly as they appear
in the data. Return nothing else — no explanation, no prose, just the JSON array.
If there are no active medications, return an empty array: []}
\end{quote}

\noindent
For strategies B, C, and D, which use dashes to indicate missing dosage fields, an
additional sentence is appended clarifying that a dash represents missing data, not
inactive status. No chain-of-thought prompting or few-shot examples are used; the
prompt is designed to elicit clean structured output directly. Temperature is set to
0.0 for all runs.

\subsection{Evaluation}
\label{sec:evaluation}

We parse the model output as a JSON array of strings and compute precision,
recall, and F1 against the ground truth medication name set using \textit{exact
string match}. A predicted medication counts as correct only if it appears
verbatim in the ground truth set. This is intentionally strict: name variations
such as ``Metformin 500 MG Oral Tablet'' versus ``metformin 500mg'' count as a
miss. F1 under exact match therefore represents a lower bound on real-world
performance, where fuzzy matching would recover some of these near-misses.
If the model output cannot be parsed as a JSON array, we record
\texttt{parse\_failed=True} and assign precision = recall = F1 = 0.

We run 200 patients $\times$ 4 strategies $\times$ 5 models = 4,000 total
inference runs.

For statistical comparisons, we use the Wilcoxon signed-rank test
\citep{wilcoxon1945individual}, a non-parametric test appropriate for the
non-normal, bounded F1 distributions we observe. For within-model strategy
comparisons (Strategy A vs.\ Strategy C), p-values are Bonferroni-corrected
for four simultaneous tests per model ($k=4$). Effect sizes are reported as
$r = |Z| / \sqrt{N}$, where $N$ is the number of patient pairs, following
\citet{field2009discovering}. We use the thresholds $r \geq 0.1$ (small),
$r \geq 0.3$ (medium), and $r \geq 0.5$ (large).
